\section{Conclusion and Future Work}
\label{sec:conclusion}

This paper proposes a concrete discipline for access-layer comparisons: derive expected read results from the declared task, check mutual equivalence and post-write state before a timing is admitted, define each configured treatment by a predeclared policy, and report fixed-demand latency separately from capacity and matched utilization. The selection policy is reproducible evidence about the treatments it fixes, not observed practitioner behavior, and the common-SQL raw-path experiment is a compound sensitivity contrast, not a library-intrinsic overhead estimate.

In the pinned Express.js experiment, the largest observed gap occurs on the tested deep fetch: $7.01\times$ on PostgreSQL and $5.12\times$ on MySQL, narrowing, in a separate ten-run campaign, to $1.71\times$ and $2.03\times$ under common SQL through raw facilities. At stable 50\% and most 70\% capacity targets, p99 values form a much smaller band than at the high-load point. At the stated evidential bar only one treatment reverses durably across the two tested backend stacks, and it is the one whose MySQL driver differs, so the experiment records limited, treatment-specific non-transfer without identifying the DBMS as its source.

The evidence record separates completed checks from unperformed validation: the oracle, preflight,
second same-host campaign, checklist, and retrospective audit make the method inspectable, while the
audit has one rater, every adapter is single-author, and both campaigns share one virtualized host.
Result-blind packets are provided instead of a claim of independent validation.

Future work should estimate inter-rater agreement for the external audit, repeat the reduced RQ2
matrix on a physically independent host, and extend the workload across datasets, topologies, and
cold-cache regimes. The first is prepared rather than proposed: a rater-extensible audit record and a
committed agreement scorer ship with the artifact.
